\documentclass[11pt]{article}
\usepackage{amsmath}
\begin{document}

\section*{Step 44: Holographic MMI Consequence}

Using the standard sign convention
\[
I_3(A:B:C)=S(A)+S(B)+S(C)+S(ABC)-S(AB)-S(AC)-S(BC),
\]
the holographic monogamy inequality is \(I_3\le 0\). This is the sign
convention equivalent to \(I(A:B)+I(A:C)\le I(A:BC)\).

On the Step-42 contracted random tensor-network carrier with
\(A=L0\), \(B=L1\), and \(C=R0\), all tested rows satisfy the inequality.
The largest holographic value is
\[
\max I_3 = -0.0255459067644\le 0.
\]
The separate min-cut entropy computation also gives \(I_3\le 0\) at every
tested bond dimension.

The non-geometric four-party GHZ control has
\[
I_3^{\rm GHZ}=0.69314718056>0,
\]
so it violates MMI. Therefore MMI is not a universal quantum-state identity;
it is a falsifiable forbidden-region signature of the geometric RT/min-cut
structure on this finite carrier.

\end{document}
